\documentclass{article}
\usepackage{parskip}
\usepackage{tabularx}
\usepackage{booktabs}
\usepackage{parskip}

\setlength{\parindent}{0pt}

\usepackage[margin=1in]{geometry}
\title{Problem Statement and Goals\\\progname}

\author{\authname}

\date{}

\input{../Comments.text}
\input{../Common.text}

\begin{document}

\maketitle

\newpage{}

\begin{table}[hp]
\caption{Revision History} \label{TblRevisionHistory}
\begin{tabularx}{\textwidth}{llX}
\toprule
\textbf{Date} & \textbf{Developer(s)} & \textbf{Change}\\
\midrule
2026/09/24 & Mohammad Bilal & Completed sections 1.1, 4, \& reflection \\
2026/09/25 & Jake Finlay & Completed section 1.4, Reflection \\
2026/09/26 & Ahsan Muzammil & Section 2 Goals, Section 3 Stretch Goals, Reflection (Ahsan) \\
2026/09/26 & Jackson Beach & Completed section 1.2 and Reflection (Jackson) \\
2026/09/27 & Mohammad Bilal & Updated reflection \\
2026/09/27 & Ethan McMehen & Added content to 1.3 Stakeholders, Reflection \\
\bottomrule
\end{tabularx}
\end{table}

\newpage{}

\section{Problem Statement}

\subsection{Problem}

Paediatric ophthalmology appointments can be stressful and unfamiliar for children.
Procedures such as vision testing, eye drops, and slit-lamp examinations may cause anxiety, particularly when children do not know what to expect.
This anxiety can make it difficult to cooperate during examinations, potentially affecting the quality and completeness of clinical assessments.
For example, a study of 298 children undergoing cycloplegic eye examinations found that 26\% experienced significant distress before receiving eye drops,
while uncooperative children were also more anxious and required more time to complete the examination \cite{ChiaEtAl2015}.

Children and their families therefore need adequate preparation for common ophthalmology procedures before attending an appointment.
Without appropriate preparation, unfamiliar clinical experiences may contribute to increased anxiety, reduced cooperation, and longer or more difficult appointments.
This project addresses the need for an accessible, child-friendly way to help children become more familiar and comfortable with the paediatric ophthalmology experience.

\subsection{Inputs and Outputs}

\textbf{Inputs:} Touchscreen selections made by children or parents to
choose preparation modules and navigate procedure walkthroughs. As well, responses
to prompts from the game, also via the touchscreen, but possibly through the use of the device's
microphone and camera too.

\textbf{Outputs:} Visual and audio content that introduces ophthalmology procedures
and facilitates participation, which includes illustrations, animations, and verbal
instruction. The application also provides feedback on activity responses. A matching
response receives confirmation, while one that does not receives
encouragement and a hint to try again. Completing an activity shows an
on-screen completion indicator, like a completion message or animation.

\subsection{Stakeholders}
    \subsubsection{Primary Stakeholders}
        \textbf{Children}

        The main audience for this interactive application is children with 
        ophthalmology visits. The application will introduce common procedures 
        to the children through age-appropriate interactive games to reduce anxiety 
        and increase preparedness for their visits. The children will be directly 
        involved with the evaluation of the application's usability, feasibility 
        and effectiveness in reducing anxiety and promoting cooperation.

        \textbf{Parents}

        Parents will be using this application together with their children due to the 
        young age range of the children. Parents will have the option and ability to 
        play with their children to assist in the introduction of the common 
        procedures through the games. Parents can be the primary source of evaluation 
        feedback of applications usability and feasibility within home environments 
        together with their children.

    \subsubsection{Secondary Stakeholders}
        \textbf{Ophthalmologists}

        Ophthalmologists will be involved with the game development, requirements 
        elicitation and design of the interactive games as experts of the clinical 
        procedures children will experience. The application can be used by 
        ophthalmologists to evaluate the preparedness of the children through the 
        successful completion of the games. 

        \textbf{Child-life Specialists}

        Child-life specialists will be involved with the game development, 
        requirements elicitation and design as the experts of how children like to 
        play and prepare for upcoming visits to reduce fear and anxiety. The 
        application can also be used by child-life specialists in a play setting 
        if the child-life specialist is present with the children, giving valuable 
        feedback on usability within the home environment.

    \subsubsection{Tertiary Stakeholders}

        \textbf{Child Psychology + Development Professors}

        Child pyschology and child development professors from McMaster University 
        can help to assist in the requirements and design development of the interactive 
        application because of their expertise in child development at various age ranges.

        \textbf{Regulatory Bodies}
        
        Regulatory bodies will be ensuring that the interactive application is working in 
        an ethical manner within ophthalmology environments and clinics. This is to ensure 
        the technology will be safe for children and their families within home environments.

        \textbf{Other Healthcare Providers}

        Other healthcare providers may benefit from the results of this interactive 
        application and use information or research provided by the results to develop 
        similar technologies within their paediatric healthcare applications. An example 
        could include a family doctor preparing a young child for a health check-up.

\subsection{Environment}

\subsubsection{Software}

The application will run on iOS and Android operating systems. The application will be distributed through the Apple App Store and Google Play Store.

\subsubsection{Hardware}

The application will run on consumer smartphones and tablets. An internet connection will be required to download the application.

\section{Goals}
\begin{enumerate}
    \item \textbf{Reduce Children's Anxiety}: 
    The application will reduce children's anxiety toward common ophthalmology procedures 
    (i.e., vision testing, eye drops, slit-lamp examinations, etc.) by helping them understand 
    what to expect before their appointment. Success will be measured by comparing children's 
    reported anxiety before and after using the application, as well as comparing anxiety 
    levels between children who use the application for preparation and those who do not. 
    Child appropriate anxiety measures, along with parent or clinician observations where 
    appropriate, can be used for these comparisons.

    \item \textbf{Improve Children's Cooperation During Examinations:}
    The application will improve children's cooperation during
    ophthalmology examinations, allowing practitioners to complete
    assessments with fewer interruptions or repeated attempts.
    Success will be measured through clinician rated cooperation,
    successful completion of examination procedures, and the
    number of attempts required to complete selected procedures.
    Additionally, children will complete an age-appropriate
    satisfaction survey, using visual rating scales such as
    smiley faces, to assess how comfortable and prepared they
    felt during their examination.

    \item \textbf{Maintain Engagement Across Age Groups}: 
    The application will provide age appropriate and engaging activities that encourage 
    children across the target age range (i.e., ages 3 to 7) to actively participate in 
    their preparation. Success will be measured through activity completion rates, usability 
    feedback, and child reported enjoyment and engagement, with results compared across 
    age groups.

    \item \textbf{Enable Effective At Home Preparation}: 
    The application will enable children and their families to prepare for ophthalmology 
    appointments independently before arriving at the clinic. Success will be measured 
    through usability and feasibility testing with children and parents, including successful 
    completion of preparation activities and parent feedback on whether the application 
    helped their child feel prepared for the appointment.

    \item \textbf{Improve Clinical Efficiency}: 
    The application will reduce difficulties and time associated with completing paediatric 
    ophthalmology examinations by preparing children before their appointment. Success will 
    be measured by comparing examination duration, repeated procedure attempts, and 
    practitioner rated examination difficulty between children who use the application and 
    children who do not.
\end{enumerate}
\section{Stretch Goals}

\begin{enumerate}
    \item \textbf{Personalize Preparation}: The application will personalize preparation 
    activities according to the procedures expected during each child's upcoming appointment, 
    ensuring that children focus on relevant examinations rather than completing the same activities 
    regardless of their visit. Success will be measured by evaluating whether personalized preparation 
    improves engagement, anxiety, or cooperation compared with a standardized preparation experience.

    \item \textbf{Provide Immersive Procedure Simulations}: The application will use augmented reality 
    to provide more realistic simulations of common ophthalmology procedures, giving children additional 
    exposure to the clinical experience before their appointment. Success will be measured by comparing 
    anxiety, preparedness, and cooperation between children using the standard application and those using 
    the augmented reality simulations.
\end{enumerate}
\section{Custom Documents (CDs)}

The custom documents chosen for the project are a \textbf{Client Interview Report} and a \textbf{Usability Report}.

The Client Interview Report will document the interviews conducted with users and stakeholders.
It will summarize the main points and feedback from the interviews and how they helped guide the project.

The Usability Report will document the results of usability testing and the feedback received from users.
This will include the testing process, results, and discussion.

\newpage{}
\setlength{\parskip}{0.8em}
\setlength{\parindent}{0pt}

\section{Appendix --- Reflection}

\subsection*{Mohammad Bilal}

\subsubsection*{What went well while writing this deliverable?}

One thing that went well was that the problem itself was relatively clear to the team and I.
The project description gave a good starting point for identifying the main issue and
explaining why it matters. Meeting with the supervisor was also helpful in clarifying the project's
scope and better understanding the overall vision for the project. The team was then able to connect 
the problem to the clinical context and identify the main users and stakeholders involved.

Another thing that went well was that the team was all on the same page. There weren't any
major disagreements, and we had a shared understanding of the problem we were trying to
solve and our vision for how to proceed.

\subsubsection*{What pain points did you experience during this deliverable, and how did you resolve them?}

One of the main challenges was deciding how much detail to include without making the problem statement
too focused on the implementation. It was easy to start describing specific app features instead of
focusing on the actual problem being addressed. The team and I resolved this by going back to the
purpose of the deliverable and keeping the problem statement focused on the clinical experience,
the needs of the users, and the overall problem.

Another challenge was the timeline for the first deliverable. At first, we were too focused on picking
a project and supervisor, so we did not get started on this deliverable as early as we could have.
Once that was done, the first deliverable was due very soon after, so I had to stay on top of my work
and make sure I was completing my tasks on time. We were able to manage this through our use of GitHub
and good organization. By properly splitting up the work into issues and evenly distributing the tasks,
everyone knew what they were responsible for and what needed to be completed. After that, the rest
of the process went smoothly.

\subsection*{Jake Finlay}

\subsubsection*{What went well while writing this deliverable?}

What made this deliverable straightforward for me was the amount of discussion
we had already had with our TA and our stakeholders before writing anything
down. By the time we sat down to write, the problem was well understood by
everyone. Gathering the information for the deliverable was mostly a matter
of writing down what we had already agreed on rather than working it out from
scratch. The other thing that helped was that we had broken the deliverable into
issues on GitHub ahead of time and assigned them to team members. That meant
there was never any ambiguity about who owned what or what needed to be picked
up next. This allowed me to move straight into writing instead of spending time
coordinating.

\subsubsection*{What pain points did you experience during this deliverable,
and how did you resolve them?}

The main pain point for me was collaborating through GitHub. Our team had not
yet agreed on conventions for pull requests and commit messages, so everyone
took a different approach. This made reviews take much longer than they needed
to. The same issue appeared in the deliverables themselves. Each person's
writing was strong on its own, but details such as formatting varied between
team members and had to be aligned. We resolved this by discussing these
differences in pull request comments and in our team chat. These discussions
were necessary because our development plan was still being written and we had
nothing to follow at the time of writing.

\subsection*{Ahsan Muzammil}

\subsubsection*{What went well while writing this deliverable?}

One aspect that went well for me while writing this deliverable was developing the core goals for the project. The project description gave me a clear starting point for 
understanding the main problems the application is intended to address. From there, I was able to identify the outcomes that I thought were most important and connect them 
back to the problem statement. This helped me develop goals around reducing children's anxiety, improving cooperation during examinations, maintaining engagement across 
different age groups, enabling at-home preparation, and improving clinical efficiency.

I also found it helpful to think about how each goal could actually be measured. Rather than only stating what I wanted the application to accomplish, I considered potential 
ways to determine whether the goal had been achieved, such as measuring anxiety levels, examination duration, repeated procedure attempts, activity completion, and user 
feedback. This made me realize that a goal needs to be specific enough that its success can be evaluated.

\subsubsection*{What pain points did you experience during this deliverable, and how did you resolve them?}

One pain point I experienced while working on this deliverable was determining the appropriate level of detail for the project goals. Initially, I found it difficult to 
distinguish between a goal and a feature of the application. For example, I considered ideas such as gamification, augmented reality, and AI personalization as potential goals 
because they were important parts of how I imagined the application working. However, I realized that these describe how the application could be implemented rather than 
what it should ultimately accomplish.

I resolved this by going back to the problem statement and asking myself what outcome each proposed goal was intended to achieve. This helped me shift my focus from 
specific features to broader outcomes, such as reducing anxiety, improving cooperation, and improving clinical efficiency. I also made each goal more measurable by 
considering how I could determine whether the application actually achieved it. This process helped me create goals that were more closely connected to the problem we 
were trying to solve and gave me a better understanding of how project goals should be defined.

Another pain point I experienced during this deliverable was maintaining consistency with our version control practices. Since we had not yet completed the coding 
standards and workflow sections of our Development Plan, we had not agreed on conventions such as how commits should be named or what structure should be used for 
pull requests. As different team members began contributing to the deliverable, there were differences in the notation and structure they used. This made it more 
difficult for me to follow a consistent format when making my own commits and pull requests.

I resolved this by looking at the existing commits and pull requests from the team and trying to follow the patterns that were already being used where possible. 
However, this also showed me that relying on existing practices is not an ideal long term solution when there is no agreed-upon standard. Going forward, I think 
establishing the coding and version control conventions earlier in the development process would make collaboration easier and reduce inconsistencies. In particular, 
agreeing on a commit naming convention and pull request template before everyone begins contributing would give the team a clear standard to follow.

\subsection*{Ethan McMehen}

\subsubsection*{What went well while writing this deliverable?}

I specifically wrote the stakeholders section of the document and I found that 
explicitly documenting the stakeholders and their relationships in the tier list 
really helped to identify who the application is being developed for. By writing 
the descriptions of each stakeholder and their relation to the problem statement, 
it really helped me to prepare for requirements elicitation and who to approach 
for what kind of information. I also found that writing a problem statement is 
fantastic too because it helps to focus on our specific problem we are aiming to 
solve. Our project description by our supervisor was solid, but writing our own 
version helps us to have a definition in an engineering context that we can refer 
to throughout the project. Overall, I learned that by having explicit definitions 
and descriptions for topics or information that seems relatively straight forward 
is still important to have a unified mission.

\subsubsection*{What pain points did you experience during this deliverable, and how did you resolve them?}

My pain points that I experienced during this deliverable have to be creating 
the right descriptions for the stakeholders and their relationship to the 
application that we are going to build for the capstone. While the stakeholders 
of children, parents, ophthalmologists etc. is easy to name, actually coming 
down to their effect and purpose in the project required some thinking. To 
resolve this pain point, I broke down their role into three parts: are they 
using the app, are they evaluating the app effectiveness, and do they care about 
the safety and ability for the app to reduce stress. I used these points to 
devise descriptions about each stakeholder relative to their involvment in each 
of these three parts. This helped me reinforce my learning that breaking down 
the problem or pain point into smaller pieces is a great approach for software 
documentation.

\subsection*{Jackson Beach}

\subsubsection*{What went well while writing this deliverable?}
The section I had in this deliverable was relatively straightforward to do, and thus I was able to accomplish it quickly and
address my larger tasks in the development plan.

\subsubsection*{What pain points did you experience during this deliverable, and how did you resolve them?}
Personally, my section of the document was not the most complex. However, it was difficult in the sense that it
can be tricky to describe inputs and outputs of the system with an appropriate level of abstraction. On the one hand,
if you go too low level, it reads as implementation details, which does not provide enough flexibility at this stage of
the project. On the other hand, if you keep it too high-level and vague, then it can be confusing to the reader, and also
make them less actionable when you go to develop the system. I resolved this by trying to strike a balance between the two,
mainly by focusing on ``what'' will be inputted by the user and outputted by the system during the normal game flow, and also
by asking for feedback from my teammates on the level of abstraction I was using was too little or too much.

\subsection*{Group Answer(s)}

\subsubsection*{How did you and your team adjust the scope of your goals to ensure they are suitable for a Capstone project?}

When defining our goals, one of the main ways we adjusted the scope was by distinguishing between what the application needs to accomplish and additional features that may be 
more difficult to implement within the time and resources available for a Capstone project. For example, augmented reality and AI based personalization were proposed as potential 
additions to the application. After discussing the scope and considering the complexity involved in implementing and evaluating these technologies, we decided to treat them as 
stretch goals rather than core requirements.

We also ensured that our core goals focused on achievable outcomes rather than attempting to address every aspect of the paediatric ophthalmology experience. Instead of having 
a broad goal such as improving the overall ophthalmology experience for children, we focused on specific outcomes such as reducing anxiety, improving cooperation during 
examinations, maintaining engagement across age groups, enabling at home preparation, and improving clinical efficiency. This provided the project with sufficient complexity 
for a senior design project while keeping the goals specific enough to be evaluated.

Overall, adjusting the scope helped us recognize that a Capstone project does not need to include every interesting idea we develop. Our core goals should represent the 
essential value of the proposed solution, while more technically complex ideas can be treated as stretch goals if time and resources allow.

\bibliographystyle{plain}
\bibliography{../../refs/References}


\end{document}
